\begin{thebibliography}{99}
\bibitem{repo}
V. Reshetnikov, \emph{ProveIt: Polylogarithms manuscript}, source snapshot
\qpin. In particular the main manuscript, integration chapter,
research-status chapter, and full-Cayley obstruction continuation.
\url{https://github.com/VladimirReshetnikov/ProveIt/tree/446dca9bd65c20e409d1d2561d4ffad408103e2a/Analysis/Polylogarithms/docs/manuscript}.

\bibitem{arithmetic}
\emph{Arithmetic spectral jets}, ProveIt research continuation,
section ``Periodic weights and several distinct arithmetic poles'', same pin.
\qfile{Analysis/Polylogarithms/docs/reports/stieltjes-harmonic-resonance/continuations/arithmetic-spectral-jets/sections/04-arithmetic.tex}.
The inspected section credits the classical simple-pole anomaly mechanism.

\bibitem{classification}
\emph{Herglotz rational classification}, ProveIt research report,
section ``Complete formal classification at rational arguments'', same pin.
\qfile{Analysis/Polylogarithms/docs/reports/herglotz-rational-classification/sections/04-classification.tex}.
This is the source of the editorial status correction, not a new
classification claimed in this article.

\bibitem{intake}
ProveIt, \emph{Incoming reports}, intake instructions and archive catalogue,
same pin. \qfile{docs/incoming/README.md} and \qfile{docs/incoming/}.
Only archive metadata, not archive interiors, was inspected.

\bibitem{hurwitz}
NIST Digital Library of Mathematical Functions, \S25.11,
\emph{Hurwitz zeta function}.
\url{https://dlmf.nist.gov/25.11}. Consulted 10 October 2026.

\bibitem{polylog}
NIST Digital Library of Mathematical Functions, \S25.12,
\emph{Polylogarithms}.
\url{https://dlmf.nist.gov/25.12}. Consulted 10 October 2026.

\bibitem{bessel}
NIST Digital Library of Mathematical Functions, \S10.32,
\emph{Integral representations of modified Bessel functions}, especially
(10.32.9) and (10.32.10).
\url{https://dlmf.nist.gov/10.32}. Consulted 10 October 2026.

\bibitem{halfK}
NIST Digital Library of Mathematical Functions, \S10.39,
\emph{Relations to other functions}, especially the half-integer base case;
the higher finite formula used here is also derived by recurrence.
\url{https://dlmf.nist.gov/10.39}. Consulted 10 October 2026.

\bibitem{theta}
NIST Digital Library of Mathematical Functions, \S20.7,
\emph{Theta-function identities and transformations}.
\url{https://dlmf.nist.gov/20.7}. Consulted 10 October 2026.

\bibitem{gamma}
NIST Digital Library of Mathematical Functions, \S5.5,
\emph{Functional relations}, and \S5.15, \emph{Polygamma functions}.
\url{https://dlmf.nist.gov/5.5}; \url{https://dlmf.nist.gov/5.15}.
Consulted 10 October 2026.

\bibitem{dowker}
J. S. Dowker, \emph{Calculation of the multiplicative anomaly},
arXiv:1412.0549, 2014; corrected version 3, 24 September 2023.
\url{https://arxiv.org/abs/1412.0549}.
Cited for the classical subject and scope; the precise defect and sign
convention used in this article are proved independently.
\end{thebibliography}
